\hnsection{CONTINUOUS GROUPS AND INFINITESIMAL TRANSFORMATIONS}

Suddenly, in the autumn of 1873, Lie again took up the study of transformation groups and obtained decisive results. Insofar as it is possible to follow his line of thought from a few letters to A. Mayer written in 1873--1874 ({[}3{]}, vol.~5, pp.~584--608), he started from a ``continuous group'' of transformations on n variables

\[
x _ {i} ^ {\prime} = f _ {i} (x _ {1}, \dots , x _ {n}, a _ {1}, \dots , a _ {r}) \quad (1 \leqslant i \leqslant n)\tag{4}
\]

depending effectively† on r parameters \(a_{1},\ldots,a_{r}\) ; he observed that, if the transformation (4) is the identity for the values \(a_{1}^{0},\ldots,a_{r}^{0}\) of the parameters,‡ then the first order Taylor expansions of the \(x_{i}\) :

\[
\begin{array}{r l} f _ {i} (x _ {1}, \dots , x _ {n}, a _ {1} ^ {0} + z _ {1}, \dots , a _ {r} ^ {0} + z _ {r}) & \\ = x _ {i} + \sum_ {k = 1} ^ {r} z _ {k} X _ {k i} (x _ {1}, \dots , x _ {n}) + \dots & (1 \leqslant i \leqslant n) \end{array} \tag {5}
\]

give a ``generic'' infinitely small transformation depending linearly on the r parameters \(z_{j}\)

\[
d x _ {i} = \left(\sum_ {k = 1} ^ {r} z _ {k} \mathrm{X} _ {k i} (x _ {1}, \dots , x _ {n})\right) d t \quad (1 \leqslant i \leqslant n).\tag{6}
\]

Proceeding as in his memoir with Klein, Lie integrated the differential system

\[
\frac {d \xi_ {1}}{\sum_ {k} z _ {k} X _ {k 1} (\xi_ {1} , \dots , \xi_ {n})} = \dots = \frac {d \xi_ {n}}{\sum_ {k} z _ {k} X _ {k n} (\xi_ {1} , \dots , \xi_ {n})} = d t,\tag{7}
\]

which gave him, for each point \((z_{1},\ldots,z_{r})\) a one-parameter group

\[
t \mapsto x _ {i} ^ {\prime} = g _ {i} (x _ {1}, \dots , x _ {n}, z _ {1}, \dots , z _ {r}, t) \quad (1 \leqslant i \leqslant n)\tag{8}
\]

such that \(g_{i}(x_{1},\ldots,x_{n},z_{1},\ldots,z_{r},0)=x_{i}\) for all i. He showed by an ingenious method, using the fact that the transformations (4) form a set which is stable under composition, that the one-parameter group (8) is a subgroup of the given group {[}3 d{]}. The new idea, the key to the whole theory, is to take the Taylor expansions of the functions (4) to the second order. The progress of his argument was quite confused and heuristic ({[}3 d{]} and {[}3{]}, vol.~5, pp.~600--601); it can be presented as follows. For sufficiently small \(z_{r}\) , let t=1 in (8); thus new parameters \(z_{1},\ldots,z_{r}\) are obtained for the transformations of the group (this is in fact the first appearance of the ``canonical parameters''). Then by definition, using (7),

\[
\frac {\partial g _ {i}}{\partial t} = \sum_ {k} z _ {k} \mathrm{X} _ {k i} (x _ {1} ^ {\prime}, \dots , x _ {n} ^ {\prime})
\]

whence

\[
\begin{array}{r l} \frac {\partial^ {2} g _ {i}}{\partial t ^ {2}} & = \sum_ {k, j} z _ {k} \frac {\partial \mathrm{X} _ {k i}}{\partial x _ {j}} (x _ {1} ^ {\prime}, \ldots , x _ {n} ^ {\prime}) \frac {\partial x _ {j} ^ {\prime}}{\partial t} \\ & = \sum_ {k, j} z _ {k} \frac {\partial \mathrm{X} _ {k i}}{\partial x _ {j}} (x _ {1} ^ {\prime}, \ldots , x _ {n} ^ {\prime}) \left(\sum_ {n} z _ {h} \mathrm{X} _ {h j} (x _ {1} ^ {\prime}, \ldots , x _ {n} ^ {\prime})\right) \end{array}
\]

which gives

\[
\begin{array}{r} x _ {i} ^ {\prime} = x _ {i} + \Big (\sum_ {k} z _ {k} \mathrm{X} _ {k i} (x _ {1}, \ldots , x _ {n}) \Big) t \\ + \frac {1}{2} \Big (\sum_ {k, n, j} z _ {k} z _ {h} \frac {\partial \mathrm{X} _ {k i}}{\partial x _ {j}} (x _ {1}, \ldots , x _ {n}) \mathrm{X} _ {h j} (x _ {1}, \ldots , x _ {n}) \Big) t ^ {2} + \dots , \end{array}
\]

whence, for \(t = 1\) , the Taylor expansions in the parameters \(z_{j}\)

\[
x _ {i} ^ {\prime} = x _ {i} + \left(\sum_ {j} z _ {k} \mathrm{X} _ {k i}\right) + \frac {1}{2} \left(\sum_ {k, h, j} z _ {k} z _ {h} \mathrm{X} _ {h j} \frac {\partial \mathrm{X} _ {k i}}{\partial x _ {j}}\right) + \dots (1 \leqslant i \leqslant n).\tag{9}
\]

We abridge these relations to \(x' = \mathbf{G}(x, z)\) relating the vectors

\[
x = (x _ {1}, \dots , x _ {n}), \quad x ^ {\prime} = (x _ {1} ^ {\prime}, \dots , x _ {n} ^ {\prime}), \quad z = (z _ {1}, \dots , z _ {r});
\]

the fundamental stability property of the set of these transformations under composition can be written

\[
\mathrm{G} (\mathrm{G} (x, u), v) = \mathrm{G} (x, \mathrm{H} (u, v))\tag{10}
\]

where \(\mathbf{H} = (\mathbf{H}_1, \ldots, \mathbf{H}_r)\) is independent of \(x\) ; it is immediate that \(\mathrm{H}(u, 0) = u\) and \(\mathrm{H}(0, v) = v\) , whence the expansions

\[
\mathrm{H} _ {i} (u, v) = u _ {i} + v _ {i} + \frac {1}{2} \sum_ {h, k} c _ {i k h} u _ {h} v _ {k} + \dots ,\tag{11}
\]

where the terms omitted are not linear in u or v. Transforming (10), using (9) and (11) and then comparing the terms in \(u_{h}v_{k}\) of the two sides. Lie obtained the relations

\[
\sum_ {j = 1} ^ {n} \left(\mathrm{X} _ {h j} \frac {\partial \mathrm{X} _ {k i}}{\partial x _ {j}} - \mathrm{X} _ {k j} \frac {\partial \mathrm{X} _ {h i}}{\partial x _ {j}}\right) = \sum_ {l = 1} ^ {r} c _ {l h k} \mathrm{X} _ {l i} (1 \leqslant h, k \leqslant r, 1 \leqslant i \leqslant n). \tag{12}
\]

His experience with the theory of partial differential equations led him to write these conditions in a simpler form: following the pattern of (1), he associated with each of the infinitely small transformations obtained by setting \(z_{k} = 1\) , \(z_{h} = 0\) for \(h \neq k\) in (6), the differential operator

\[
\mathrm{A} _ {k} (f) = \sum_ {i = 1} ^ {n} \mathrm{X} _ {k i} \frac {\partial f}{\partial x _ {i}},\tag{13}
\]

and rewrote conditions (12) in the form

\[
[ \mathrm{A} _ {h}, \mathrm{A} _ {k} ] = \sum_ {l} c _ {l h k} \mathrm{A} _ {l},\tag{14}
\]

the corner stone of his theory. Until then he used the terms ``infinitely small transformation'' and ``infinitesimal transformation'' indifferently (e.g.~{[}3 e{]}); the simplicity of relations (14) led him to call the operator (13) the ``symbol'' of the infinitesimal transformation \(dx_{i} = X_{ki}dt\) ( \(1 \leqslant i \leqslant n\) ) {[}3 d{]} and very soon he called the operator (13) itself the ``infinitesimal transformation'' ({[}3 d{]} and {[}3{]}, vol.~5, p.~589).

He then became aware of the close links which united the theory of ``continuous groups'' with his earlier research on contact transformations and partial differential equations. This bringing together filled him with enthusiasm: ``My earlier works were as it were all ready there waiting to found the new theory of transformation groups'' he wrote to Mayer in 1874 ({[}3{]}, vol.~5, p.~586).

In the following years, Lie pursued the study of transformation groups. Besides the general theorems summarized below (§ III), he obtained some more special results: the determination of the transformation groups of the line and plane, the subgroups of low codimension in the projective groups, the groups in at most 6 parameters, etc. He did not abandon differential equations for long. In fact, it seems that, for him, the theory of transformation groups was a tool for integrating differential equations, where the transformation group played a role analogous to that of the Galois group of an algebraic equation.†

We note that this research led him equally to the introduction of certain transformation sets in an infinity of parameters which he called ``infinite continuous groups'' \(^{‡}\) ; he reserved the name ``finite continuous groups'' for transformation groups in a finite number of parameters of type (4) above.

